%% PHANTOM: physics-informed harmonic-attention transit vetting
%% Prepared for Astronomy and Computing (Elsevier).  The manuscript is written
%% against elsarticle; switching to mnras.cls requires only the preamble.
\documentclass[preprint,12pt,authoryear]{elsarticle}

\usepackage{amsmath}
\usepackage{amssymb}
\usepackage{graphicx}
\usepackage{booktabs}
\usepackage{siunitx}
\usepackage{subcaption}
\usepackage{xcolor}
\usepackage{tabularx}
\usepackage[hidelinks]{hyperref}

\journal{Astronomy and Computing}

\newcommand{\phantomnet}{\textsc{Phantom}}
\newcommand{\astronet}{\textsc{AstroNet}}
\newcommand{\exominer}{\textsc{ExoMiner}}

\begin{document}

\begin{frontmatter}

\title{Harmonic attention and a differentiable transit model for
       automated vetting of exoplanet candidates}

\author[a]{Anonymous Author}
\address[a]{Affiliation}

\begin{abstract}
%% Filled in once the final numbers are in; see RESULTS_PLACEHOLDER markers.
\end{abstract}

\begin{keyword}
exoplanets \sep transit photometry \sep light curves \sep deep learning \sep
conformal prediction \sep Kepler \sep TESS
\end{keyword}

\end{frontmatter}

%% ===================================================================
\section{Introduction}
\label{sec:intro}

The transit method has produced the large majority of known exoplanets, but the
detections it yields are not planets: they are \emph{threshold-crossing events}
(TCEs), periodic dimmings that survive a matched-filter search of a photometric
time series. Most TCEs are not planetary. They arise from eclipsing binary
stars, from the light of an eclipsing binary blended into the photometric
aperture of an unrelated target, from stellar variability and starspots, and
from instrumental systematics that happen to be periodic. The Kepler mission
alone produced $\sim$$3.4\times10^{4}$ TCEs in its final data release, of which
only a minority correspond to planets \citep{Thompson2018}. Deciding which is
which --- \emph{vetting} --- is the bottleneck between a photometric survey and
a planet catalogue.

Vetting has traditionally been performed by a cascade of physically motivated
diagnostic tests. If odd-numbered and even-numbered transits have different
depths, the signal is probably an eclipsing binary detected at half its true
period. If a shallower dimming appears near the opposite orbital phase, the
companion is self-luminous and therefore stellar. If the flux-weighted centroid
of the target shifts during the event, the light is being removed from a
neighbouring source rather than from the target. If the profile is V-shaped
rather than flat-bottomed, the eclipse is probably grazing. These tests are
fast, interpretable, and were the basis of the Kepler Robovetter
\citep{Thompson2018}. They are also brittle: each is a hand-tuned threshold on a
statistic that degrades rapidly as the signal-to-noise ratio falls, and each
fails in a characteristic way. The odd-even test, in particular, is defeated by
exactly the situation it is meant to catch, because a binary detected at its
half-period alias presents both eclipses as ``primary'' events.

Deep learning changed the picture. \citet{Shallue2018} showed that a
one-dimensional convolutional network applied to two fixed-length
representations of the phase-folded light curve --- a whole-orbit ``global''
view and a zoomed ``local'' view --- outperformed the hand-built cascade, and
used it to recover two previously missed planets. Subsequent work added input
channels rather than changing the reasoning: \citet{Ansdell2018} supplied
centroid time series and stellar parameters, and \exominer{}
\citep{Valizadegan2022} added a branch per diagnostic test and validated 301 new
planets. The same recipe was carried to TESS \citep{Yu2019, Osborn2020}, and
more recently attention has been introduced over the same two views, together
with post-hoc temperature scaling, by \citet{Islam2026}.

Two concurrent studies bear directly on the present work. \citet{Priyanshu2026}
give the first application of conformal prediction to transit \emph{detection},
using a self-supervised stellar world model to find single transits in raw flux
and reporting empirical coverage of conformal prediction sets.
\citet{Kopparapu2026} examine cross-instrument prioritisation over all
train/test combinations of Kepler and TESS, and find that a model trained on one
mission transfers poorly to the other, recommending joint training instead. We
return to both in Section~\ref{sec:discussion}: our conformal contribution
targets a different guarantee (false discovery rate over a selected list, rather
than coverage of a prediction set) on a different task (vetting events of known
period, rather than detecting single transits), and our cross-mission experiment
is deliberately zero-shot, which is the setting \citet{Kopparapu2026} identify as
difficult.

Across this line of work the network is given the folded light curve
\emph{at the period reported by the search} and asked to classify what it sees.
That framing embeds an assumption which the classical tests explicitly do not
make: that the reported period is the true period. It is precisely the failure
of that assumption --- the half-period alias --- which produces the dominant
astrophysical false positive. A network shown only the fold at $P$ can learn
proxies for the alias, but it cannot examine the alternative hypothesis
directly, because the evidence that distinguishes $P$ from $2P$ lives in a fold
the network is never shown.

A second gap concerns what the network outputs. Vetting produces a ranked list
that consumes finite follow-up resources, so the operationally meaningful
question is not ``what is the score of this candidate'' but ``if I observe
everything on this list, what fraction of my time is wasted''. A softmax output
does not answer that question: neural network probabilities are systematically
miscalibrated \citep{Guo2017}, and calibration on one data set carries no
guarantee on another. Nor do these classifiers estimate the physical parameters
--- period, depth, duration --- that the survey exists to measure; those are
taken from the upstream search and passed through unchanged.

This paper addresses both gaps with a single architecture, \phantomnet{}
(Physics-informed Harmonic-ATtention Object Model), and makes four
contributions.

\begin{enumerate}
\item \textbf{Harmonic-hypothesis attention.} The light curve is rendered under
      three competing period hypotheses --- $P/2$, $P$ and $2P$ --- alongside
      the odd, even, secondary and centroid channels. A transformer encoder
      attends over the resulting token set, and an explicit contrast head
      compares the pooled evidence under each hypothesis. The network can
      therefore test the alias hypothesis in the same way a human vetter does,
      rather than inferring it indirectly (Section~\ref{sec:harmonic}).
\item \textbf{A differentiable transit decoder.} A five-parameter analytic
      transit model --- depth, half-duration, phase offset, limb darkening and
      an ingress-softness term that interpolates continuously between a
      flat-bottomed and a V-shaped profile --- is rendered onto the local view
      and fitted by reconstruction, with no parameter supervision. The residual
      becomes a classification feature, and the bottleneck yields transit
      parameters as a by-product (Section~\ref{sec:decoder}).
\item \textbf{Distribution-free candidate lists.} Conformal $p$-values computed
      against held-out known non-planets, combined with the
      Benjamini--Hochberg procedure, produce a candidate list whose false
      discovery rate is controlled at a chosen level under an exchangeability
      assumption alone --- with no assumption that the network is well
      calibrated or even accurate (Section~\ref{sec:conformal}). This controls
      the quantity a follow-up programme actually spends its time on, namely the
      fraction of the selected list that is not planetary, and is complementary
      to the coverage guarantee reported for single-transit detection by
      \citet{Priyanshu2026}.
\item \textbf{An evaluation protocol that does not leak.} Multiple TCEs on the
      same star share a light curve, stellar properties and systematics.
      Splitting at the level of the TCE, as previous work has done, places such
      siblings on both sides of the train/test boundary. We report results under
      a star-level split as our primary protocol and the TCE-level split for
      comparability (Section~\ref{sec:splits}).
\end{enumerate}

%% ===================================================================
\section{Data}
\label{sec:data}

\subsection{The Kepler DR24 threshold-crossing event set}

We use the Kepler Q1--Q17 DR24 TCE catalogue hosted at the NASA Exoplanet
Archive, which is the benchmark introduced by \citet{Shallue2018} and reused by
\citet{Ansdell2018}. Of its 20{,}367 TCEs, 15{,}737 carry a disposition from the
Autovetter training set \citep{McCauliff2015}:

\begin{itemize}
\item \textbf{PC} (planet candidate, 3{,}600 events) --- the positive class;
\item \textbf{AFP} (astrophysical false positive, 9{,}596 events) --- eclipsing
      binaries, blends and other genuine astrophysical signals that are not
      planetary transits of the target;
\item \textbf{NTP} (non-transiting phenomenon, 2{,}541 events) --- stellar
      variability and instrumental artefacts.
\end{itemize}

These 15{,}737 events lie on 9{,}865 distinct stars. We treat PC as the positive
class and $\{$AFP, NTP$\}$ as the negative class, following the convention of
the prior work we compare against.

Light curves are the long-cadence ($29.4$~min) PDCSAP photometry from the Kepler
pipeline, retrieved for every quarter in which the target was observed. We
retain cadences whose quality flag contains none of the bits indicating attitude
tweaks, safe modes, coarse or Earth point, desaturation events, manual
exclusions or detector anomalies.

\subsection{The TESS cross-mission set}
\label{sec:tessdata}

To test whether the learned representation is specific to Kepler we assemble a
zero-shot evaluation set from the TESS Objects of Interest catalogue
\citep{Guerrero2021}. We keep TOIs whose TFOPWG disposition is unambiguous ---
confirmed or known planets (CP, KP) as positives, false positives and false
alarms (FP, FA) as negatives --- and which have a published epoch, period and
duration. Photometry is the two-minute-cadence SPOC light curve, limited to the
first four sectors available for each target; only cadences with a zero quality
flag are kept, since the TESS quality bits differ from Kepler's.

%% ===================================================================
\section{From photometry to fixed-length views}
\label{sec:preprocessing}

\subsection{Detrending}

Stellar variability and residual instrumental trends exceed a planetary transit
in amplitude by orders of magnitude, so each light curve is flattened before
folding. Following \citet{Vanderburg2014} and \citet{Shallue2018} we divide each
contiguous segment --- data separated by gaps longer than $0.75$~d are treated
as separate segments --- by a cubic basis spline fitted with iterative
$3\sigma$ outlier rejection.

Two details matter. First, the knot spacing is not fixed: it is selected per
target from a logarithmic grid spanning $0.5$ to $20$~d by minimising the
Bayesian Information Criterion summed over all segments, so that a rapidly
rotating star and a photometrically quiet one are treated differently. Second,
the cadences within $1.5$ transit durations of \emph{any} TCE on the target are
excluded from the spline fit, though the fitted spline is still evaluated at
those cadences. Without this mask the spline partially absorbs the very signal
being measured, biasing depths low.

The centroid time series are detrended in the same way and with the same knot
spacing, so that only excursions coincident with an event survive.

\subsection{Views}
\label{sec:views}

For each TCE, cadences falling within $1.5$ durations of a \emph{different} TCE
on the same star are removed, and the remainder is folded on the reported
ephemeris. Nine fixed-length views are then rendered by median-binning
(Table~\ref{tab:views}). Each is centred on its median and divided by the
absolute value of its minimum, so that the deepest bin sits at $-1$; the divisor
is retained as a scalar, which is what allows an absolute depth to be recovered
from the normalised view in Section~\ref{sec:decoder}.

\begin{table}[t]
\centering
\caption{The nine input views. The period hypothesis column indicates which
         trial period the view is folded on; these labels are supplied to the
         network as embeddings.}
\label{tab:views}
\begin{tabular}{llcl}
\toprule
View & Bins & Hypothesis & Purpose \\
\midrule
Global       & 2001 & $P$   & whole-orbit morphology \\
Local        & 201  & $P$   & transit shape \\
Odd          & 201  & $P$   & depth on odd epochs \\
Even         & 201  & $P$   & depth on even epochs \\
Secondary    & 201  & $P$   & occultation by a luminous companion \\
Harmonic     & 201  & $P/2$ & is the true period half of the reported one? \\
Harmonic     & 2001 & $2P$  & does the fold separate into unequal eclipses? \\
Centroid (global) & 2001 & $P$ & photocentre motion, whole orbit \\
Centroid (local)  & 201  & $P$ & photocentre motion during the event \\
\bottomrule
\end{tabular}
\end{table}

The two harmonic views are the novel channels. Folding at $2P$ maps alternating
eclipses of a binary onto \emph{distinct} phases, so a depth difference that the
odd-even test must infer from two interleaved subsets becomes directly visible
as two separated dimmings of unequal depth. Folding at $P/2$ tests the opposite
hypothesis: if the true period is half the reported one, that fold sharpens the
event instead of smearing it.

The local, odd, even and secondary views span $\pm4$ transit durations. The
secondary view is centred by re-folding about the epoch of the deepest
out-of-primary bin, rather than by windowing inside the primary fold, so that
the window remains fully populated when the secondary lies near phase $\pm0.5$.

\subsection{Scalar features}

Alongside the views the network receives scalar features in four groups:
\emph{transit} (period, duration, depth, model signal-to-noise, multiple event
statistic, number of transits, impact parameter, radius and semi-major axis
ratios), \emph{stellar} (effective temperature, surface gravity, metallicity,
inferred planet radius), \emph{diagnostic} (the pipeline's odd-even and centroid
offset statistics and robust fit statistic), and \emph{derived} --- eleven
quantities measured here from the detrended light curve, including the odd-even
depth asymmetry, the secondary depth ratio, the selected knot spacing, the
out-of-transit scatter, the normalisation scale, and a red-noise proxy formed
from the ratio of the scatter of six-hour bins to the white-noise expectation.
The derived group is the only one computable identically for both missions and
is therefore the configuration used for the cross-mission experiment.
Features spanning orders of magnitude are log-compressed, missing values are
imputed to the training-split median, and all are standardised with
training-split statistics.

%% ===================================================================
\section{The \phantomnet{} architecture}
\label{sec:model}

\phantomnet{} has three parts: a per-view convolutional encoder, a harmonic
attention block over the resulting tokens, and a differentiable transit decoder.

\subsection{View encoders}

Each view is encoded by a stack of pre-activation residual blocks with
GroupNorm, interleaved with strided max-pooling --- five stages for the 2001-bin
views and three for the 201-bin views. The final feature map is reduced by
concatenated mean and max pooling, capturing respectively the average level and
the depth of the deepest excursion, and projected to a token of dimension
$d=256$. Views of equal length share an architecture but not weights, since a
centroid view and a flux view carry different meaning.

\subsection{Harmonic-hypothesis attention}
\label{sec:harmonic}

Each token receives two additive embeddings: a view-type embedding and a
\emph{period-hypothesis} embedding drawn from $\{P/2,\,P,\,2P\}$. A learnable
classification token is prepended and the sequence is processed by a
pre-norm transformer encoder \citep{Vaswani2017} of three layers and eight
heads.

Attention alone leaves the harmonic comparison implicit, so we add an explicit
contrast. Writing $h_{P/2}$, $h_{P}$ and $h_{2P}$ for the mean of the encoded
tokens under each hypothesis, the contrast vector
\begin{equation}
c = \big[\,h_{P}\ \|\ h_{P}-h_{P/2}\ \|\ h_{P}-h_{2P}\ \|\ |h_{P}-h_{2P}|\,\big]
\label{eq:contrast}
\end{equation}
is passed through a single hidden layer and concatenated into the classifier
input. The difference terms state directly whether the evidence prefers the
reported period over its harmonics; the absolute-value term makes the comparison
insensitive to the sign of the discrepancy, which is what the eclipsing-binary
hypothesis requires.

\subsection{Differentiable transit decoder}
\label{sec:decoder}

Let $x$ index the local-view grid in units of transit durations from
mid-transit. From the encoded local-view token an affine head predicts five
bounded parameters: depth $\delta>0$, half-duration $T\in(0.05,2)$, phase offset
$x_0\in(-0.5,0.5)$, limb-darkening strength $u\in(0,1)$ and ingress softness
$s\in(0.02,1)$. With $z = |x-x_0|/T$ the model light curve is
\begin{equation}
\hat{f}(x) = -\,\delta\;
  \underbrace{\sigma\!\left(\frac{1-z}{s}\right)}_{\text{in-transit indicator}}\;
  \underbrace{\Big[1-u\big(1-\sqrt{\max(1-z^{2},0)}\big)\Big]}_{\text{limb darkening}},
\label{eq:render}
\end{equation}
where $\sigma$ is the logistic function. Equation~\ref{eq:render} is analytic and
differentiable in all five parameters, so it can be trained by backpropagation
through a reconstruction loss against the observed local view.

The softness $s$ is the physically interesting parameter. As $s\to0$ the
indicator approaches a step and the profile becomes flat-bottomed, the signature
of a small opaque body crossing the stellar disc. As $s$ grows the ingress and
egress ramps broaden until the profile is V-shaped, the signature of a grazing
or blended stellar eclipse. A single scalar therefore expresses the U-versus-V
distinction that the classical shape test tries to capture with a ratio of
widths, but it is fitted jointly with depth, duration and limb darkening rather
than measured on a noisy curve in isolation.

Two quantities derived from the fit enter the classifier: the root-mean-square
reconstruction residual and its maximum absolute value. A genuine transit is
well described by Equation~\ref{eq:render}; a stellar eclipse or an instrumental
artefact is not, and the residual records that mismatch.

Critically, the five parameters are \emph{not} regressed against catalogue
values. Supervising them would teach the network to reproduce the pipeline fit
against which it is later compared, making any parameter-recovery result
circular. They are identified by reconstruction alone.

\subsection{Objective}

The training objective is
\begin{equation}
\mathcal{L} = \mathcal{L}_{\mathrm{BCE}}(\hat{y}, y)
  \;+\; \lambda\,\big\| \hat{f} - f_{\mathrm{local}} \big\|_2^2 ,
\label{eq:loss}
\end{equation}
with $\lambda=5$. Networks are optimised with AdamW \citep{Loshchilov2019} under
a one-cycle schedule, batch size 128, gradient clipping at norm 5, and early
stopping on validation average precision. Views are randomly mirrored about
mid-transit during training, which is a valid symmetry of a transit profile.

%% ===================================================================
\section{Evaluation protocol}
\label{sec:protocol}

\subsection{Splits}
\label{sec:splits}

Our primary protocol splits $80/10/10$ into training, validation and test sets
\emph{by star}: every TCE on a given target falls in exactly one fold. Multiple
TCEs on one star share a light curve, the same detrending solution, the same
stellar parameters and the same instrumental systematics, so a TCE-level split
places near-duplicates on both sides of the boundary and inflates test scores.
Because previous work reports the TCE-level split, we also report it, so that
the comparison to published numbers is like-for-like and the size of the leakage
is itself measurable.

\subsection{Metrics}

We report the area under the ROC curve, average precision (the area under the
precision-recall curve), and precision at fixed recall. The last is the
operational quantity: a vetting pipeline is run at high recall so as not to
discard planets, and is then judged by how much follow-up effort the surviving
false positives consume. All figures are averaged over five random seeds and
reported with their standard deviation; ensembles average the per-seed
probabilities.

\subsection{Distribution-free candidate selection}
\label{sec:conformal}

Let $s$ denote the model score. Treating ``this event is not a planet'' as the
null hypothesis, and given $n$ scores $s_1,\dots,s_n$ of \emph{known
non-planets} from the held-out validation split, the conformal $p$-value of a
test event with score $s$ is
\begin{equation}
p = \frac{1 + \#\{i : s_i \ge s\}}{n+1} .
\label{eq:pvalue}
\end{equation}
Applying the Benjamini--Hochberg procedure \citep{Benjamini1995} at level $q$ to
these $p$-values yields a candidate list whose false discovery rate is
controlled at $q$ \citep{Bates2023}. The guarantee is finite-sample and
distribution-free: it requires only that calibration and test events be
exchangeable, and holds regardless of whether the network is well specified or
well calibrated. For comparison we also report the expected calibration error
before and after temperature scaling \citep{Guo2017}.

\subsection{Transit parameter recovery}

The decoder bottleneck is converted to physical units by
$\delta_{\mathrm{ppm}} = \delta\,\kappa\times10^{6}$, where $\kappa$ is the
normalisation scale divided out of the local view, and
$T_{\mathrm{dur}} = 2T\,T_{\mathrm{cat}}$, where $T_{\mathrm{cat}}$ is the
catalogue duration defining the grid. Uncertainties are the standard deviation
across ensemble members \citep{Lakshminarayanan2017}.

These estimates are compared against transit parameters from the NASA Exoplanet
Archive's composite planetary parameters table --- that is, against published
follow-up analyses rather than against the Kepler pipeline whose fit is being
assessed. The comparison is therefore between two independent estimates of the
same physical quantity: the pipeline value that seeded the search, and the value
recovered here.

%% ===================================================================
\section{Results}
\label{sec:results}

%% RESULTS_PLACEHOLDER

%% ===================================================================
\section{Discussion}
\label{sec:discussion}

%% DISCUSSION_PLACEHOLDER

%% ===================================================================
\section{Conclusions}
\label{sec:conclusions}

%% CONCLUSIONS_PLACEHOLDER

\section*{Data availability}

All photometry and catalogues used here are public: Kepler and TESS light curves
from the Mikulski Archive for Space Telescopes, and the DR24 TCE, cumulative KOI,
TOI and composite planetary parameter tables from the NASA Exoplanet Archive.
Code to reproduce every figure and table is released with this paper.

\bibliographystyle{elsarticle-harv}
\bibliography{references}

\end{document}
